%%%%%%%%%%%%%%%%%%%%%%%%%%%%%%%%%%%%%%%%%%%%%%%%%%%%%%%%%%%%%%%%%%%%%%%%%
%
%	LaTeX File for PhD Thesis - Background and Related Work
%
%%%%%%%%%%%%%%%%%%%%%%%%%%%%%%%%%%%%%%%%%%%%%%%%%%%%%%%%%%%%%%%%%%%%%%%%%

\chapter{Background and Related Work}
\label{chap:background}

\section{Introduction}
% TODO: Write introduction positioning the background chapter
% This chapter surveys the foundational concepts and prior research relevant to this thesis.

\section{Related Work in Automatic Modulation Classification}
\label{sec:related_work}

This section reviews the progression of methods for Automatic Modulation Classification (AMC), which serves as the primary validation task for the contributions in this dissertation. The review begins with the transition from traditional feature-based techniques to the application of Convolutional Neural Networks (CNNs). It then discusses the subsequent adoption of Transformer-based architectures, which demonstrated improved performance by capturing long-range dependencies in signal data. Finally, it covers the emergence of self-supervised learning paradigms aimed at reducing the reliance on large labeled datasets. This survey of the literature reveals the absence of a unified, modular framework that integrates these capabilities—a gap this dissertation aims to address.

\subsection{Early Deep Learning for Automatic Modulation Classification}

Initial research in Automatic Modulation Classification (AMC) relied on likelihood-based methods, which required precise channel state information, or feature-based methods using hand-crafted statistics like higher-order cumulants and cyclostationary features~\cite{swami2000hierarchical,gardner1994cyclostationarity,nandi1995automatic}. The introduction of deep learning, particularly Convolutional Neural Networks (CNNs)~\cite{o2016convolutional,peng2017modulation}, marked a significant advancement by enabling the automatic learning of discriminative features from raw signal data. Early work demonstrated that CNNs applied to 2D signal representations, such as constellation diagrams, could achieve high accuracy, albeit with large training datasets~\cite{peng2018modulation}. Subsequent research focused on architectural refinements, such as specialized kernels and multi-representation frameworks, to better capture the spatio-temporal characteristics of wireless signals~\cite{huynh2020mcnet}.

\subsection{Transformer-Based Architectures in AMC}

While effective for learning local patterns, the limited receptive fields of CNNs restrict their ability to model long-range dependencies in signal sequences, a notable drawback in low-SNR environments. Transformer architectures~\cite{vaswani2017attention} were introduced to address this limitation. Their self-attention mechanisms proved highly effective at modeling the global context of a signal, leading to superior performance over CNN-based approaches, particularly in noisy conditions~\cite{cai2022signal}. This led to further research integrating convolutional embeddings to ground the attention mechanism~\cite{kong2023transformer} and developing hybrid architectures that combined Transformers with LSTMs to better model temporal dynamics~\cite{qu2024enhancing}. The successful application of Vision Transformers (ViTs)~\cite{dosovitskiy2020image} to constellation diagrams further validated that attention-based models, when properly adapted, achieve superior generalization with minimal data~\cite{}.

\subsection{Self-Supervised Learning for Data Efficiency}

A primary limitation of the aforementioned supervised models is their dependency on large labeled datasets, which are often costly and difficult to acquire in the wireless domain. To mitigate this, research has increasingly focused on self-supervised learning (SSL) to learn effective representations from abundant unlabeled data. SSL paradigms from other domains, such as masked modeling from natural language processing~\cite{devlin2019bert,he2022masked} and contrastive learning from computer vision~\cite{chen2020simple,he2020momentum}, have been successfully adapted for wireless signals. Our work, DenoMAE~\cite{faysal2025denomae}, contributes to this area by extending masked autoencoding to a multimodal framework that treats noise as an explicit and informative modality, a concept detailed further in Chapter~\ref{chap:denomae}. While these advances address specific challenges, they highlight the need for a unified framework that combines these capabilities economically.

\subsubsection{Multisensory Masked Modeling}

Masking large portions of the input and training the model to reconstruct or predict missing elements has become a popular self-supervised approach, especially within autoencoder-based methods. For example, BERT \cite{devlin2018bert} introduced masked language modeling by masking words in sentences and training the model to predict them, a concept adapted for vision with Masked Autoencoders (MAE), where image patches are masked to encourage the model to learn robust visual representations \cite{he2021mae}.

In multimodal contexts, masked modeling extends to tasks combining text, vision, and audio, promoting cross-modal understanding. For instance, Vision-Language Pretraining (VLP) models \cite{lu2019vilbert, chen2020uniter, wang2022image} utilize masked language and vision modeling to jointly learn from masked image regions and text tokens. Additionally, image modality-centered models like MultiMAE \cite{bachmann2022multimae}, 4M \cite{mizrahi20244m}, and 4M21 \cite{bachmann20244m} have demonstrated strong performance by pretraining on multiple image modalities and enhancing inference capabilities across diverse inputs. These multimodal masking methods have demonstrated significant gains in self-supervised learning, advancing model performance in tasks that benefit from multimodal reasoning.

\subsection{Challenges in Wireless Communication}

Building foundation models for wireless signals requires addressing fundamental challenges that distinguish this domain from vision and language:

\begin{itemize}
    \item \textbf{SNR Imbalance and Condition-Specific Performance:} Existing models overlook critical distinctions between low and high SNR conditions\cite{}, resulting in uneven classification performance. Low SNR signals obscured by noise require fundamentally different feature extraction strategies compared to high SNR signals\cite{}. Traditional monolithic models cannot adequately handle this performance imbalance\cite{}. We propose to address this through DenoMAE's hybrid generative-discriminative learning that captures both global and local patterns, combined with an improved plug-and-play discriminative classifier\cite{}, designed to handle varying SNR conditions robustly, ensuring consistent performance across the full SNR spectrum.
    
    \item \textbf{Data Scarcity and Labeling Costs:} Acquiring high-quality labeled RF signal data demands specialized equipment and expertise, making large-scale labeled dataset collection prohibitively expensive\cite{}. Unlike vision or language domains where crowdsourcing enables massive labeled datasets\cite{}, wireless signals require domain experts and controlled RF environments for annotation. Foundation models must leverage self-supervised learning\cite{} on abundant unlabeled signals to dramatically reduce labeled data requirements, achieving high accuracy with minimal supervision.
    
    \item \textbf{Lack of Semantic Grounding and Interpretability:} Images depict recognizable objects; text conveys explicit meaning. Wireless signals encode information through abstract statistical properties (symbol distributions, spectral characteristics, temporal patterns) that lack inherent semantic content\cite{}. Current deep learning AMC systems provide no interpretability—they are black boxes incapable of explaining decisions or enabling human-understandable interaction\cite{}. Foundation models must bridge this semantic gap, converting quantitative signal features into structured natural language representations that enable both LLM reasoning and human comprehension.
    
    \item \textbf{Signal Corruption and Noise Variability:} Unlike additive noise in images, wireless signals experience channel-dependent corruption including fading, multipath interference, and time-varying distortion\cite{}. Signal-to-noise ratios can range from -10 dB to 30+ dB within the same system, requiring representations robust across this spectrum\cite{}. Traditional approaches treat noise as pure corruption, but it actually contains information about channel physics and transmission conditions\cite{} that foundation models should exploit as an informative modality.
    
    \item \textbf{Open-Set Reality and Continuous Evolution:} Wireless systems continuously encounter new modulation schemes, emerging standards (e.g., 6G), and adversarial signals\cite{}. Models must recognize when they encounter novel signal types rather than forcing everything into known categories—a capability absent in current closed-set AMC systems\cite{}. Foundation models require representations enabling not just classification of known schemes but detection and characterization of novel signals through compositional understanding and reasoning\cite{}.
    
    \item \textbf{Multimodal Representational Complexity:} The same wireless signal manifests differently across time-domain IQ samples, frequency spectra, constellation diagrams, and cyclostationary features\cite{}, with varying lengths and sampling rates. Each representation emphasizes different signal properties\cite{}. Foundation models must integrate these complementary views through multimodal processing\cite{} to build comprehensive understanding, while remaining flexible to diverse signal formats.
    
    \item \textbf{Task Diversity and Rapid Adaptation Requirements:} A wireless foundation model should support not just classification but also denoising, channel estimation, signal generation, and cross-modulation transfer\cite{} with minimal task-specific modification. Unlike vision models trained primarily for recognition\cite{}, wireless models need representations enabling both analysis and synthesis. Moreover, they must rapidly adapt to new deployment scenarios with few labeled examples, requiring efficient meta-learning capabilities\cite{} beyond standard transfer learning\cite{}.
    
    \item \textbf{Deployment Cost and Computational Constraints:} Training massive foundation models from scratch for each wireless deployment scenario is computationally prohibitive\cite{}. Real-world wireless systems require cost-effective solutions that leverage existing world knowledge—from pre-trained LLMs' reasoning capabilities to self-supervised representations learned once and transferred broadly—rather than exhaustive training across all possible channel conditions, modulation schemes, and SNR ranges.
\end{itemize}

These challenges motivated our modular, cost-effective framework: learning universal representations through multimodal masked autoencoders (addressing data scarcity and multimodal complexity), enabling training-free reasoning through LLM integration (providing interpretability and open-set capability), ensuring condition-specific robustness through modular architectures (handling SNR imbalance), and achieving rapid adaptation through meta-learning (enabling efficient knowledge transfer).

\subsection{Cross-Domain Challenges in Representation Learning}

Building foundation models requires representations that transcend specific tasks and transfer broadly. Both wireless communication and meta-learning face common representation learning challenges:

\begin{itemize}
    \item \textbf{Task-Feature Misalignment:} Standard masked autoencoding objectives (predict missing regions) learn representations optimized for reconstruction rather than discrimination or reasoning\cite{}. For wireless signals, reconstructing masked IQ samples may not capture modulation-specific features necessary for classification or open-set recognition\cite{}. The gap between self-supervised pretext tasks and downstream objectives limits transfer effectiveness\cite{}.
    
    \item \textbf{Modality Integration:} Wireless signals naturally exist in multiple modalities (time-domain, frequency-domain, constellation diagrams)\cite{}, each emphasizing different properties. Naive concatenation or late fusion fails to capture cross-modal dependencies\cite{}. Foundation models must learn joint representations where modalities inform and disambiguate each other, especially under noise.
    
    \item \textbf{Noise as Modality vs. Corruption:} Standard representation learning treats noise as corruption to overcome\cite{}. In wireless systems, noise patterns encode channel state and signal quality information crucial for adaptation\cite{}. Learning to leverage noise as an informative modality rather than merely being robust to it represents a paradigm shift requiring novel architectures.
    
    \item \textbf{Computational Efficiency:} Large-scale self-supervised learning typically requires enormous batch sizes (SimCLR needs 4096+\cite{chen2020simple}) and computational resources\cite{}. Wireless edge devices and real-time systems demand efficient alternatives\cite{}. Foundation model pre-training must balance representation quality with practical deployment constraints.
    
    \item \textbf{Semantic Grounding:} Vision models leverage natural image statistics\cite{}; language models exploit linguistic structure\cite{}. Wireless signals lack inherent semantic content—their "meaning" emerges from technical specifications and physical principles\cite{}. Grounding wireless representations in natural language descriptions (e.g., "QPSK with raised cosine filtering") could enable reasoning and zero-shot generalization through language model integration.
    
    \item \textbf{Supervised-Unsupervised Integration:} Effectively combining labeled and unlabeled data remains challenging\cite{}. Most approaches treat them as separate learning stages (pre-train then fine-tune) rather than joint objectives\cite{}. For wireless foundation models, simultaneously leveraging massive unlabeled signals and limited labeled examples through unified loss functions could improve both representation quality and task performance.
    
    \item \textbf{Open-Set Recognition Gap:} Standard representation learning assumes closed-set downstream tasks\cite{}. Foundation models must produce representations that enable not just classification of known classes but detection and characterization of novel signal types\cite{}. This requires representations capturing interpretable, compositional features rather than just discriminative boundaries.
\end{itemize}

Addressing these challenges requires rethinking representation learning architectures, objectives, and integration strategies—the core focus of this dissertation.

% \section{Foundation Models in AI}
% % TODO: GPT, BERT, CLIP, MAE, ViT

% \subsection{Natural Language Processing}
% % TODO: GPT series, BERT, transformer architectures

% \subsection{Computer Vision}
% % TODO: CLIP, MAE, ViT, self-supervised learning in vision

% \subsection{Key Principles of Foundation Models}
% % TODO: Pre-training, fine-tuning, transfer learning, scale

% \section{Wireless Communications and Signal Processing}
% % TODO: Overview of wireless communication fundamentals

\subsection{RF-Based Drone Detection and Classification}
\label{sec:drone_detection_background}

The rapid proliferation of Unmanned Aerial Vehicles (UAVs) has created urgent needs for reliable detection and classification systems. While multiple sensing modalities exist (radar, acoustics, vision, RF), this section focuses on RF-based approaches relevant to Chapter~\ref{chap:cdrf}. For comprehensive surveys of detection methodologies, signal representations, and public datasets, see~\cite{allahham2020deep, basak2021combined}.

\subsubsection{Detection Modalities}

RF-based sensing offers unique advantages: day/night operation, non-line-of-sight detection, and exploitation of persistent command-and-control links. However, it faces challenges from co-channel interference (Wi-Fi, Bluetooth) and limited performance at low SNR~\cite{nguyen2025effective, rfuav}. Other modalities include radar (long range but struggles with small targets), acoustics (low cost but noise-sensitive), vision (high accuracy but weather-dependent), and multi-sensor fusion (robust but complex)~\cite{akter2022explainable, mrabet2024machine}.

\subsubsection{RF Signal Representations}

Visual time-frequency representations dominate modern RF perception. Spectrograms via STFT expose center frequency, bandwidth, and timing patterns suitable for CNNs~\cite{rfuav}. Alternative encodings include wavelets, Wigner–Ville distributions, and FDGAF mappings~\cite{bai2020intelligent}. Empirical evidence shows spectrogram-based models outperform raw I/Q approaches under interference~\cite{fu2024radio, medaiyese2022wavelet}.

\subsubsection{Deep Learning Approaches}

CNN architectures (ResNet, specialized variants) dominate classification tasks~\cite{KILIC2022101028, he2016deep}. Object detectors like YOLO jointly localize and classify emissions~\cite{rfuav}. Recent work explores Transformers~\cite{vaswani2017attention}, sequence models (LSTM), and classical ML (XGBoost)~\cite{drone2024daiarticle}. A key challenge: models trained on clean datasets show inflated scores but poor real-world generalization~\cite{elyousseph2024robustness}.

\subsubsection{Dataset Landscape}

Public datasets include DroneRF~\cite{allahham2019dronerf}, DroneDetect~\cite{5jjj-1m32-21}, Cardinal RF, and RFUAV~\cite{rfuav}. Common limitations: narrow class diversity, restricted raw I/Q access, poorly characterized SNR, limited negative samples, and lack of integrated tooling for augmentation and evaluation. Chapter~\ref{chap:cdrf} addresses these gaps.

% \subsection{Modulation and Demodulation}
% % TODO: Basic concepts, modulation schemes

\subsection{Higher-Order Statistics and Cumulants}
\label{sec:higher_order_stats}

Higher-order statistics and cumulants provide powerful tools for signal characterization, particularly relevant to the plug-and-play LLM approach in Chapter~\ref{chap:plugandplay}. This section summarizes the key theoretical foundations.

Statistical moments quantitatively describe the shape and characteristics of a random variable's probability distribution. The $n^{\mathrm{th}}$-order moment of a random variable $x$ is:
\begin{equation}
    m_n(x) = \int x^n f(x) \, dx,
\end{equation}
where $f(x)$ is the probability density function. Central moments $m_n(x-\mu)$ and normalized moments $m_n(x/\sigma)$ provide additional characterization.

Cumulants $c_n$ offer an alternative distribution characterization. The first three correspond to mean, variance, and skewness; higher-order cumulants capture additional shape aspects. The cumulant-generating function:
\begin{equation}
    C_t(x) = \log\left(E\left[e^{tx}\right]\right) = \sum_{n=1}^{\infty} \frac{c_n t^n}{n!}.
\end{equation}

Cumulants exhibit additive behavior for independent random variables: $c_n(X+Y) = c_n(X) + c_n(Y)$. In practice, cumulants are computed from moments via relationships like:
\begin{align}
  c_{4,0} &= m_{4,0} - 3\,m_{2,0}^2,\\
  c_{4,2} &= m_{4,2} - \lvert m_{2,0}\rvert^2 - 2\,m_{2,1}^2,
\end{align}
where mixed moments $m_{q,p} = \frac{1}{N} \sum_{n=1}^{N} x[n]^q \, x[n]^{p-q}$~\cite{AZZOUZ199555}. These statistics enable effective signal characterization for modulation classification.

\section{Meta-Learning and Few-Shot Learning}
\label{sec:meta_learning_background}

Meta-learning, or "learning to learn," enables models to rapidly adapt to new tasks with minimal data. This section reviews key concepts and regularization techniques relevant to Chapter~\ref{chap:meta_task}.

\subsection{Few-Shot Learning Paradigm}

Few-Shot Learning (FSL) addresses data scarcity by learning from limited examples. The N-Way K-Shot paradigm uses K labeled examples per class for N classes, with support sets for training and query sets for evaluation. Meta-learning approaches like Prototypical Networks~\cite{proto}, MAML~\cite{maml}, and Relation Networks~\cite{relation} enable effective generalization from sparse datasets.

\subsection{Regularization in Meta-Learning}

Regularization prevents overfitting and enhances meta-generalization. Techniques include:

\textbf{Explicit Regularization (ER)} imposes regularization terms on meta-learning updates (iMAML~\cite{imaml}, MR-MAML~\cite{mrmaml}), directly constraining optimization.

\textbf{Episode Augmentation Regularization (EAR)} modifies tasks through noise or mixup~\cite{yamaguchi2023regularizing, shu2023dac}, expanding training data.

\textbf{Task Augmentation Regularization (TAR)} generates new tasks via interpolation or augmentation~\cite{yao2021meta, rajendran2020meta}. Methods like MetaMix~\cite{chen2021metamix} and ProtoDiff~\cite{du2023protodiff} expand task distributions for robust generalization.

\textbf{Adversarial Task Upsampling (ATU)} generates tasks matching true distributions via adversarial loss~\cite{wu2022adversarial}.

Recent work like Meta-Baseline~\cite{chen2021meta} bridges transfer learning and meta-learning through pre-trained extractors. However, these approaches often require large-scale pre-training or computationally intensive operations. Our Meta-Task framework (Chapter~\ref{chap:meta_task}) introduces learnable, task-specific regularization through auxiliary tasks, learning robust embeddings during meta-training without extensive pre-training.

\subsection{Channel Effects and Noise}
% TODO: AWGN, fading, SNR

\subsection{Automatic Modulation Classification}
% TODO: AMC problem definition, traditional approaches

% \section{Self-Supervised Learning}
% % TODO: Self-supervised learning paradigm

% \subsection{Contrastive Learning}
% % TODO: SimCLR, MoCo, contrastive objectives

% \subsection{Masked Autoencoders}
% % TODO: MAE in vision, BERT in NLP, masked modeling

% \subsection{Multimodal Learning}
% % TODO: Multi-modal self-supervised learning

% \section{Large Language Models for Non-Text Domains}
% % TODO: LLMs applied beyond text

% \subsection{LLMs for Vision}
% % TODO: GPT-4V, LLaVA, vision-language models

% \subsection{Time Series and Signal Processing}
% % TODO: LLMs for time series, signal processing applications

% \subsection{Prompt Engineering and In-Context Learning}
% % TODO: Prompt design, in-context learning, few-shot prompting

% \section{Meta-Learning and Few-Shot Learning}
% % TODO: Meta-learning overview

% \subsection{Episodic Training}
% % TODO: Episode construction, support/query sets

% \subsection{Model-Agnostic Meta-Learning}
% % TODO: MAML and variants

% \subsection{Meta-Regularization}
% % TODO: Regularization approaches in meta-learning

\subsection{Limitations in Few-Shot Learning}

Meta-learning provides crucial capabilities for foundation models—enabling rapid adaptation to new tasks with minimal data. However, applying meta-learning to wireless signals reveals several limitations:

\begin{itemize}
    \item \textbf{Representation Quality Bottleneck:} Most meta-learning approaches focus on adaptation mechanisms (how to fine-tune) rather than representation learning (what features to learn)\cite{}. Poor initial representations limit adaptation effectiveness regardless of meta-learning sophistication\cite{}. Standard episodic training typically learns task-agnostic embeddings that miss domain-specific structure\cite{}.
    
    \item \textbf{Episode Construction Suboptimality:} Meta-learning performance critically depends on task distribution during training\cite{}. Random episode sampling creates tasks of varying difficulty and informativeness\cite{}. Easy episodes waste computation; impossible episodes provide no signal. Optimal episode construction for wireless signals—considering SNR variation, modulation similarity, and channel diversity—remains unexplored.
    
    \item \textbf{Unlabeled Data Underutilization:} Wireless systems generate massive unlabeled signal streams\cite{}, yet meta-learning typically requires labeled support sets\cite{}. Each episode contains rich self-supervised learning opportunities within support and query examples that standard meta-learners ignore\cite{}. Leveraging unlabeled data within episodes could dramatically improve representation quality.
    
    \item \textbf{Cross-Domain Generalization Failure:} Meta-learners trained on synthetic datasets (like RadioML\cite{}) often fail on real over-the-air signals due to distributional shift\cite{}. The model learns to adapt within the training distribution but cannot generalize to different noise profiles, hardware impairments, or channel characteristics\cite{}. Foundation models must learn invariances that transfer across domains.
    
    \item \textbf{Overfitting to Task Distribution:} Models can overfit to the specific set of modulation types, SNR ranges, and channel models in meta-training, learning spurious correlations rather than fundamental signal properties\cite{}. This limits zero-shot generalization to truly novel scenarios outside the meta-training distribution.
    
    \item \textbf{Lack of Interpretability and Reasoning:} Meta-learners adapt through gradient descent or metric learning but cannot explain their decisions or reason about signal properties\cite{}. For wireless applications requiring trust and verification (military, regulatory), pure pattern-matching adaptation is insufficient\cite{}. Integration with language models could enable interpretable reasoning.
\end{itemize}

These limitations motivated our development of episodic representation learning (MetaRep), optimized episode generation (BEpiC), and task-aware regularization (Meta-Task) specifically designed for wireless foundation model development.

\subsection{Meta-Learning for Few-Shot Adaptation}

Meta-learning provides complementary capabilities for rapid adaptation, enabling models to learn from few examples by training on distributions of tasks. Approaches like Model-Agnostic Meta-Learning (MAML)~\cite{finn2017model} and Prototypical Networks~\cite{snell2017prototypical} have shown promise in wireless contexts. However, applying meta-learning to wireless signals reveals several fundamental limitations that our research addresses.

\textbf{Episode Construction and Binary Classification:} Standard meta-learning episodes randomly select classes without considering inter-class or intra-class distance, which proves suboptimal for binary classification scenarios where one class encompasses multiple diverse categories. Our BEpiC framework~\cite{} addresses this through strategic episode generation: using PCA for dimensionality reduction to identify cluster centroids, selecting representative samples closest to centroids, and employing ORB feature matching to quantify dissimilarity between classes. On Caltech256, BEpiC achieved 63.5\% accuracy in 2-way 1-shot scenarios compared to 60.5\% for regular episodes, with even greater improvements (72.8\% vs. 70.9\%) in 5-shot settings. Combining regular and BEpiC episodes (Regular+BEpiC) further enhanced performance by effectively doubling the sample diversity, demonstrating faster convergence and improved generalization.

\textbf{Task-Aware Regularization:} Meta-learning models face significant overfitting challenges when trained on small, variable datasets, tending to memorize rather than generalize to unseen tasks. Conventional regularization techniques (weight decay, dropout, noise injection) prove insufficient for meta-learning's unique requirements of both adaptation-generalization and meta-generalization. Our Meta-Task framework~\cite{} reformulates regularization as a learnable task within the FSL optimization process, introducing parameterized neural layers that dynamically learn task-specific regularization strategies.

The Task-Decoder implementation within Meta-Task employs unsupervised autoencoders to reconstruct input images from embeddings, refining hidden representations by constraining the embedding vector to accurately reconstruct original inputs. This approach proves method-agnostic, seamlessly integrating with Prototypical Networks, MAML, MetaOptNet, and P>M>F across standard benchmarks. On MiniImageNet, Task-Decoder improved Prototypical Networks from 46.6\% to 50.0\% (5-way 1-shot) and 65.4\% to 71.3\% (5-way 5-shot) on TieredImageNet, demonstrating superior meta-generalization. Critically, these improvements occurred with faster convergence, achieving higher accuracy within 20k episodes compared to baseline methods requiring 50k episodes, while maintaining minimal computational overhead through shallow decoder architectures (75K-250K parameters).

% \section{Wireless Signal Datasets and Cross-Domain Applications}
% % TODO: Dataset overview

% \subsection{DenoMAE 2.0 Dataset Specifications}
% % TODO: Describe DenoMAE 2.0 dataset (10 classes)

% \subsection{Custom Generated Wireless Datasets}
% % TODO: Describe custom 10K dataset (10 classes, 10,000 samples)

% \subsection{Drone Detection Dataset}
% % TODO: Cross-domain validation dataset (YOLO-based)

% \subsection{Dataset Characteristics and Challenges}
% % TODO: Common challenges, preprocessing needs

% \section{Summary and Positioning}
% % TODO: Summarize background and identify gaps that this thesis addresses

\subsection{Language Model Integration for Wireless Classification}

Large Language Models (LLMs) have demonstrated remarkable capabilities in reasoning, few-shot learning, and transferring knowledge across domains through in-context learning~\cite{}. Their ability to understand structured prompts and perform inference based on descriptions presents an unexplored opportunity for wireless communications: \textit{Can we bridge wireless signal understanding and language model reasoning by converting statistical signal features into natural language prompts?}

Our plug-and-play framework~\cite{rostami2025plug} provides a pioneering proof-of-concept for this vision. By computing higher-order statistics and cumulants from wireless signals and converting them into structured natural language prompts, we enable pre-trained LLMs to perform training-free, one-shot modulation classification. Using models ranging from DeepSeek-R1-Distill-Qwen (7B and 32B parameters) to OpenAI's o3-mini (200B parameters), we demonstrated that incorporating exemplar context (I+C+S prompts) substantially improves performance compared to instruction-only prompts (I+S), achieving up to 72.4\% cleaned accuracy under noisy conditions with o3-mini.

This approach validates three critical insights: (1) sophisticated decision-making need not be trained from scratch but can leverage pre-trained LLMs' reasoning; (2) the semantic gap between statistical signal properties and linguistic descriptions can be bridged through careful prompt engineering; (3) training-free LLM-based classification achieves competitive performance compared to extensively trained neural networks, drastically reducing deployment costs while providing interpretability—a critical requirement for wireless systems.
